\section{本章小结}
\label{sec:ls-normed-inner-product-summary}

范数给向量规定大小，从而可以比较误差并定义收敛。有限维空间中的范数彼此等价，向量列收敛等价于逐坐标收敛；这些空间及其子空间都是完备的。对偶范数把输入大小的限制转化为线性泛函的最大读数，商范数则衡量忽略一个子空间中的差别以后剩余的大小。

内积进一步给出角度和正交。Gram矩阵记录一组向量的两两内积，正交分解把向量拆成子空间内的部分与垂直部分；正交投影是内积诱导距离下唯一的最近点。标准正交基简化配对与坐标计算，线性泛函的内积表示和伴随则连接了上一章的代数对偶与本章的几何结构。

这些结论为下一章寻找线性作用的简单表示提供了工具。有限维中的闭性、完备性和连续性不能不加条件地推广到无限维空间。
